% 对应题号 M1-62；图与本题条件、推导、结论逐项对应。
% 图契约：印刷宽16cm；最小字号8.4pt；几何/场或曲线在左，计算关系在右。
\documentclass[tikz,border=7pt]{standalone}
\usepackage{fontspec,xeCJK,amsmath,amssymb,graphicx}
\IfFontExistsTF{Times New Roman}{\setmainfont{Times New Roman}}{\setmainfont{TeX Gyre Termes}}
\IfFontExistsTF{Heiti SC}{\setCJKmainfont[AutoFakeBold=1.5]{Heiti SC}}{\setCJKmainfont[AutoFakeBold=1.5]{FandolHei-Regular}}
\usetikzlibrary{arrows.meta,bending,calc,patterns,positioning}
\definecolor{ink}{HTML}{243745}
\definecolor{blue}{HTML}{275DAB}
\definecolor{green}{HTML}{177C63}
\definecolor{red}{HTML}{BD4935}
\definecolor{orange}{HTML}{AA6C16}
\definecolor{purple}{HTML}{715393}
\tikzset{every node/.style={font=\fontsize{9.3}{12.5}\selectfont,text=ink},
 axis/.style={-{Stealth[length=4pt]},line width=.65pt,draw=ink!65},
 vec/.style={-{Stealth[length=5pt]},line width=1pt,draw=blue},
 short/.style={-{Stealth[length=3pt]},line width=.8pt},
 guide/.style={dashed,draw=ink!40,line width=.5pt},
 card/.style={rounded corners=2pt,draw=blue!65,fill=blue!4,align=center,text width=6.6cm,inner xsep=9pt,inner ysep=7pt},
 title/.style={font=\fontsize{10.5}{14}\selectfont\bfseries,text width=15cm,align=center},
 label/.style={font=\fontsize{8.4}{11}\selectfont,align=center}}
\begin{document}
\begin{tikzpicture}[x=1cm,y=1cm]
\path[use as bounding box] (-8,-3.4) rectangle(8,3.4);
\node[title] at (0,3) {M1-62\quad 点电荷的 δ 支撑与积分区间};
\begin{scope}[shift={(-4,0)}]

\draw[axis] (-2.7,0)--(2.7,0) node[right] {$x$};
\draw[vec,red] (0,0)--(0,1.8) node[right] {$q$};\node[below] at (0,0) {$x_0$};
\draw[green,line width=2pt] (-1.7,-.7)--(1.5,-.7);\node[label,green] at (0,-1.15) {区间 $I_1$ 包含 $x_0$};
\draw[orange,line width=2pt] (.7,-1.7)--(2.3,-1.7);\node[label,orange] at (1.5,-2.1) {$I_2$ 不包含 $x_0$};
\node[label] at (0,2.4) {箭头表示集中权重，图示 $q>0$};

\end{scope}\node[card,draw=blue!65,fill=blue!4] at (3.95,1.65) {$\rho(x)=q\delta(x-x_0)$};
\node[card,draw=green!65,fill=green!4] at (3.95,0) {$\displaystyle\int_{I_1}\rho\,dx=q$};
\node[card,draw=orange!65,fill=orange!4] at (3.95,-1.65) {$\displaystyle\int_{I_2}\rho\,dx=0$\\是否包含支撑点决定积分结果。};
\end{tikzpicture}\end{document}
